\subsection{Linear Attention}
\label{subsec:linear-attn}

Linear attention methods seek to reduce the quadratic complexity of softmax attention to $O(Nd^2)$ or $O(Nd)$ by exploiting a kernel reformulation of the attention computation. The foundational insight, formalized by Katharopoulos et al.~\cite{katharopoulos2020transformers}, is that softmax attention can be decomposed via a feature map $\phi(\cdot)$ such that $\text{softmax}(QK^\top)V \approx \phi(Q)(\phi(K)^\top V)$, where the change of associativity eliminates the explicit $N \times N$ attention matrix. This reformulation further reveals a fundamental duality: linear attention computed left-to-right is mathematically equivalent to a linear recurrent neural network with state $S_t = S_{t-1} + \phi(k_t)v_t^\top$. This duality---later formalized as the Structured State Space Duality (SSD) by Dao and Gu~\cite{dao2024transformers}---unifies linear attention, state space models, and structured matrix transformations as three views of the same underlying computation, and provides the theoretical foundation for the entire family of sub-quadratic token mixers.

The evolution of linear attention can be understood as a progression along two axes: \emph{expressivity} (from simple kernel substitution to sophisticated memory update rules) and \emph{hardware efficiency} (from naive implementations to IO-aware algorithms that fully exploit GPU tensor cores). We organize this subsection accordingly, tracing the development from kernel-based methods (2020) through gated variants (2024) to delta-rule formulations (2024--2025) and hardware-efficient training algorithms.

\subsubsection{Kernel-Based Linear Attention}
\label{subsubsec:kernel-linear}

The kernel perspective on attention begins with the observation that softmax attention can be written as a kernel function. For queries $q_i$ and keys $k_j$, the softmax attention weight $\alpha_{ij} \propto \exp(q_i^\top k_j / \sqrt{d})$ defines a kernel $\kappa(q, k) = \exp(q^\top k / \sqrt{d})$. If this kernel can be decomposed as $\kappa(q, k) = \phi(q)^\top \phi(k)$ for some feature map $\phi: \mathbb{R}^d \to \mathbb{R}^D$, then the attention output can be rewritten using the associativity of matrix multiplication.

\paragraph{Linear Transformers (Katharopoulos et al., 2020).}
Katharopoulos et al.~\cite{katharopoulos2020transformers} proposed replacing the exponential kernel with a general feature map $\phi(x) = \text{elu}(x) + 1$, yielding the linear attention formulation:
\begin{equation}
\text{LinAttn}(Q, K, V)_i = \frac{\phi(Q_i)^\top \left(\sum_{j \leq i} \phi(K_j) V_j^\top\right)}{\phi(Q_i)^\top \left(\sum_{j \leq i} \phi(K_j)\right)}.
\label{eq:linear-attn}
\end{equation}
The key insight is that the numerator and denominator can be computed incrementally via recurrent state updates:
\begin{equation}
S_i = S_{i-1} + \phi(k_i) v_i^\top, \quad z_i = z_{i-1} + \phi(k_i), \quad o_i = \frac{\phi(q_i)^\top S_i}{\phi(q_i)^\top z_i},
\label{eq:linear-attn-recurrent}
\end{equation}
where $S_i \in \mathbb{R}^{D \times d_v}$ is the recurrent state and $z_i \in \mathbb{R}^D$ is the normalizer. This formulation supports two computation modes: a \emph{parallel} mode (matrix form, $O(N^2 D)$ but parallelizable) for training, and a \emph{recurrent} mode ($O(D d_v)$ per step) for inference. The recurrent mode transforms the Transformer into an RNN with constant per-step cost, eliminating the need for a KV cache entirely. However, the simple additive state update $S_i = S_{i-1} + \phi(k_i) v_i^\top$ suffers from a fundamental limitation: information accumulates monotonically and can never be forgotten or overwritten, leading to degraded performance on tasks requiring precise associative recall.

\paragraph{Performer (FAVOR+).}
Choromanski et al.~\cite{choromanski2021rethinking} proposed approximating the softmax kernel more faithfully using random feature maps. The FAVOR+ (Fast Attention Via positive Orthogonal Random features) mechanism constructs:
\begin{equation}
\phi(x) = \frac{1}{\sqrt{m}} \exp\!\left(-\frac{\|x\|^2}{2}\right) \left[\exp(\omega_1^\top x), \ldots, \exp(\omega_m^\top x)\right]^\top,
\label{eq:favor}
\end{equation}
where $\omega_1, \ldots, \omega_m$ are orthogonal random vectors drawn from a Gaussian distribution. The orthogonality reduces the variance of the kernel approximation compared to i.i.d.\ random features, and the positive exponential ensures non-negative attention weights. With $m$ random features, the complexity is $O(Nmd)$, which is linear in $N$ when $m$ is fixed. However, achieving a good approximation of softmax requires $m = \Omega(d \log d)$ features, and the approximation quality degrades for sharp attention distributions (where a few tokens receive most of the attention mass), precisely the regime that matters most for language modeling. These limitations have prevented Performer from achieving competitive quality with softmax attention on standard LLM benchmarks.

\paragraph{Linformer.}
Wang et al.~\cite{wang2020linformer} took a different approach to linearization: rather than approximating the kernel, Linformer projects the key and value matrices to a fixed lower dimension $k \ll N$ using learned projection matrices $E, F \in \mathbb{R}^{k \times N}$:
\begin{equation}
\text{Linformer}(Q, K, V) = \text{softmax}\!\left(\frac{Q (EK)^\top}{\sqrt{d}}\right) (FV),
\label{eq:linformer}
\end{equation}
yielding $O(Nkd)$ complexity. The theoretical justification relies on the Johnson--Lindenstrauss lemma: if the attention matrix is approximately low-rank (which empirical studies confirm for typical language modeling tasks), then a random projection to $k = O(\log N)$ dimensions preserves the attention distribution with high probability. However, Linformer requires the sequence length $N$ to be known at design time (since $E$ and $F$ have $N$ columns), limiting its applicability to variable-length and autoregressive settings.

\subsubsection{Taylor-Expansion Approaches}
\label{subsubsec:taylor-linear}

\paragraph{Based.}
Arora et al.~\cite{eyuboglu2024based} proposed a principled middle ground between exact softmax and crude kernel approximation. Based uses a second-order Taylor expansion of the exponential kernel:
\begin{equation}
\exp(q^\top k) \approx 1 + q^\top k + \frac{(q^\top k)^2}{2},
\label{eq:based-taylor}
\end{equation}
which yields a feature map of dimension $O(d^2)$---larger than the $O(d)$ of simple kernels but much smaller than the $O(d \log d)$ required by random features. To compensate for the limited recall capability of the linear attention component, Based combines it with a short sliding-window softmax attention over a small local window, creating a hybrid that balances global linear attention (for throughput) with local softmax attention (for recall). The authors provide a systematic analysis of the \emph{recall--throughput trade-off}: linear attention achieves high throughput but low recall on associative retrieval tasks, while softmax attention achieves high recall but low throughput. Based's hybrid design explicitly targets the Pareto frontier of this trade-off.

\subsubsection{Gated Linear Attention}
\label{subsubsec:gated-linear}

The fundamental limitation of vanilla linear attention---monotonic state accumulation without forgetting---motivated the introduction of data-dependent gating mechanisms that control how much of the previous state is retained at each time step.

\paragraph{Gated Linear Attention (GLA).}
Yang et al.~\cite{yang2024gla} introduced a gating mechanism that modulates the state update with a data-dependent forget gate:
\begin{equation}
S_t = G_t \odot S_{t-1} + k_t v_t^\top, \quad o_t = q_t^\top S_t,
\label{eq:gla}
\end{equation}
where $G_t = \sigma(W_g x_t) \in \mathbb{R}^{d_k \times d_v}$ is a matrix-valued gate that controls, for each entry of the state matrix, how much of the previous state is retained. When $G_t = \mathbf{1}$ (all ones), GLA reduces to vanilla linear attention; when $G_t$ is a scalar times the identity, it reduces to RetNet's~\cite{sun2023retentive} exponential decay. The matrix-valued gate provides fine-grained control over which dimensions of the state are preserved and which are forgotten, enabling the model to selectively retain relevant information while discarding noise.

A key practical contribution is the \textsc{FlashLinearAttention} algorithm, which enables hardware-efficient training through a chunk-wise parallel scan. The sequence is divided into chunks of size $C$; within each chunk, the recurrence is unrolled into a matrix multiplication (exploiting GPU tensor cores), while across chunks, the state is propagated via the recurrence. This achieves $O(NC^2 + NC d_k d_v / C)$ complexity, which for optimal chunk size $C = O(\sqrt{d_k d_v})$ yields $O(N d_k d_v)$---linear in sequence length and fully hardware-efficient.

\paragraph{HGRN and HGRN2.}
Qin et al.~\cite{qin2023hgrn} introduced the Hierarchically Gated Recurrent Neural Network (HGRN), which imposes a \emph{layer-wise hierarchy} on the gating mechanism: shallow layers use gates with small lower bounds (enabling fast forgetting for local patterns), while deep layers use gates with large lower bounds (enabling slow forgetting for long-range dependencies). The lower bound for layer $l$ is set to $b_l = l/L$, providing a principled multi-scale temporal structure without additional hyperparameters. HGRN2~\cite{qin2024hgrn2} extends this framework with \emph{state expansion} via outer products: the state update becomes $S_t = G_t \odot S_{t-1} + k_t \otimes v_t$, where $S_t \in \mathbb{R}^{d_k \times d_v}$ is a matrix-valued state. This increases the state capacity from $O(d)$ to $O(d_k \times d_v)$ without additional parameters (since the outer product is parameter-free), and the resulting formulation can be interpreted as a gated linear attention mechanism, enabling chunk-wise parallel training via the same algorithms used for GLA.

\subsubsection{Delta Rule Linear Attention}
\label{subsubsec:delta-linear}

While gating addresses the forgetting problem, it does not solve the \emph{overwriting} problem: when a new key--value pair $(k_t, v_t)$ is written to the state, it does not remove the old value previously associated with key $k_t$. This leads to interference in associative memory tasks where the same key may be associated with different values at different time steps.

\paragraph{DeltaNet.}
Yang et al.~\cite{yang2024deltanet} drew inspiration from the delta rule in classical associative memory to propose a ``write-after-erase'' state update:
\begin{equation}
S_t = S_{t-1} - \beta_t k_t (k_t^\top S_{t-1}) + \beta_t k_t v_t^\top = (I - \beta_t k_t k_t^\top) S_{t-1} + \beta_t v_t k_t^\top,
\label{eq:deltanet}
\end{equation}
where $\beta_t \in (0, 1)$ is a learned step size. The first term explicitly erases the old value associated with key $k_t$ by projecting it out of the state, and the second term writes the new value $v_t$. This ``erase-then-write'' mechanism enables precise associative recall: after writing $(k, v)$ to the state, querying with $k$ retrieves exactly $v$ (up to the approximation introduced by $\beta_t < 1$). A remarkable theoretical connection is that DeltaNet's update rule is mathematically equivalent to one step of online gradient descent on the linear regression loss $\|S k_t - v_t\|^2$, establishing a deep link between linear attention and test-time training (Section~\ref{subsec:ttm}).

Parallelization of the delta rule is non-trivial because the erase operation $(I - \beta_t k_t k_t^\top)$ introduces a multiplicative dependence on the state. Yang et al.\ solve this through a compact \emph{Householder WY representation} that reformulates the sequential updates as structured matrix operations amenable to chunk-wise parallel computation, achieving training efficiency comparable to GLA.

\subsubsection{Hardware-Efficient Training and Inference}
\label{subsubsec:hw-linear}

The practical viability of linear attention depends critically on hardware-efficient implementations that exploit the memory hierarchy of modern GPUs.

\paragraph{Lightning Attention-2.}
Qin et al.~\cite{qin2024lightning} proposed a block-wise decomposition of causal linear attention that separates the computation into \emph{intra-block} and \emph{inter-block} components:
\begin{equation}
O_b = \underbrace{\text{CausalAttn}(Q_b, K_b, V_b)}_{\text{intra-block}} + \underbrace{\phi(Q_b) \cdot S_{b-1}}_{\text{inter-block}}, \quad S_b = S_{b-1} + \phi(K_b)^\top V_b,
\label{eq:lightning}
\end{equation}
where $b$ indexes blocks of size $C$. The intra-block computation uses standard causal attention (a small $C \times C$ matrix), while the inter-block computation propagates a compressed state $S_b \in \mathbb{R}^{D \times d_v}$ across blocks. This decomposition achieves constant memory consumption (independent of sequence length) and constant training speed per token, earning the characterization as a ``free lunch'' for unlimited sequence lengths. Lightning Attention-2 underpins the MiniMax-01 model~\cite{minimax2025}, which deploys it at 456B-parameter scale with 4M-token context.

\paragraph{TransNormerLLM.}
Qin et al.~\cite{qin2024transnormerllm} demonstrated that a carefully designed linear attention architecture can surpass softmax-based Transformers of comparable size on standard language modeling benchmarks. TransNormerLLM combines gated linear attention with linearized relative positional encoding with decay (LRPE-d), which introduces an exponential decay factor $\lambda^{i-j}$ into the linear attention kernel: $\alpha_{ij} = \phi(q_i)^\top \phi(k_j) \cdot \lambda^{i-j}$. This decay provides a soft recency bias without the hard cutoff of a sliding window. At 7B parameters, TransNormerLLM achieves lower perplexity than Llama on multiple benchmarks, providing the first evidence that linear attention can be competitive with softmax at production-relevant scales.

\subsubsection{Comparison and Analysis}

Table~\ref{tab:linear-attn-comparison} summarizes the key characteristics of linear attention methods along the dimensions of feature map, state size, training complexity, and recall capability.

\begin{table}[t]
\centering
\caption{Comparison of linear attention methods. $N$: sequence length, $d$: model dimension, $D$: feature dimension, $d_s = d_k \times d_v$: state size.}
\label{tab:linear-attn-comparison}
\small
\begin{tabular}{lccccc}
\toprule
\textbf{Method} & \textbf{Year} & \textbf{Feature Map / Gate} & \textbf{State Size} & \textbf{Training} & \textbf{Recall} \\
\midrule
Linear Transformer & 2020 & $\text{elu}(x)+1$ & $D \times d_v$ & $O(NDd_v)$ & Low \\
Performer          & 2020 & Random (FAVOR+) & $m \times d_v$ & $O(Nmd_v)$ & Low \\
Linformer          & 2020 & Learned projection & $k \times d_v$ & $O(Nkd)$ & Medium \\
Based              & 2024 & Taylor (2nd order) & $d^2 \times d_v$ & $O(Nd^2 d_v)$ & Medium \\
GLA                & 2024 & Data-dependent gate & $d_k \times d_v$ & $O(Nd_k d_v)$ & Medium \\
HGRN2              & 2024 & Hierarchical gate + outer product & $d_k \times d_v$ & $O(Nd_k d_v)$ & Medium \\
DeltaNet           & 2024 & Delta rule (erase+write) & $d_k \times d_v$ & $O(Nd_k d_v)$ & High \\
Lightning Attn-2   & 2024 & Block-wise decomposition & $D \times d_v$ & $O(ND d_v)$ & Low--Med \\
TransNormerLLM     & 2024 & Gate + LRPE-d decay & $D \times d_v$ & $O(NDd_v)$ & Medium \\
\bottomrule
\end{tabular}
\end{table}

The evolution from simple kernel substitution (Linear Transformer, Performer) through gated variants (GLA, HGRN2) to delta-rule formulations (DeltaNet) represents a principled trajectory toward increasingly expressive linear-complexity token mixers. Each step addresses a specific limitation of its predecessor: gating solves the monotonic accumulation problem of vanilla linear attention, and the delta rule solves the interference problem of gated linear attention. The SSD framework~\cite{dao2024transformers} provides a unifying theoretical lens, proving that all these variants---along with SSMs (Section~\ref{subsec:ssm})---are instances of the same underlying computation viewed through different mathematical lenses.

Despite this progress, linear attention methods face a fundamental expressivity ceiling: the fixed-size state $S \in \mathbb{R}^{d_k \times d_v}$ can store at most $O(d_k d_v)$ bits of information about the input history, regardless of the sequence length. This \emph{state bottleneck} manifests as degraded performance on tasks requiring precise retrieval from long contexts (e.g., needle-in-a-haystack, multi-hop reasoning). The dominant strategy for addressing this limitation is hybrid architectures (Section~\ref{subsec:hybrid}) that interleave a small number of full softmax attention layers (for precise recall) with linear attention layers (for efficient long-range processing). The independent convergence of multiple production models on a ${\sim}$7:1 linear-to-softmax ratio (Jamba, MiniMax-01) suggests a natural balance point, though whether this ratio is universal or task-dependent remains an open question.
